% appendixA.tex — Appendix A: Partitioned Matrices and Kronecker Products

\chapter{Partitioned Matrices and Kronecker Products}
\label{app:matrices}

This appendix collects some useful facts about partitioned matrices and
Kronecker products that are used throughout the text.

%% ── A.1 Addition and Multiplication ──────────────────────────────
\section{Addition and Multiplication of Partitioned Matrices}
\label{sec:A.1}

A matrix can be partitioned into submatrices (blocks) by drawing horizontal
and vertical lines between its rows and columns.  For example, the
$m \times n$ matrix $\vec{A}$ can be written as
\begin{equation}
  \vec{A} =
  \begin{bmatrix}
    \vec{A}_{11} & \vec{A}_{12} \\
    \vec{A}_{21} & \vec{A}_{22}
  \end{bmatrix},
  \label{eq:A.1}
\end{equation}
where $\vec{A}_{11}$ is $m_1 \times n_1$, $\vec{A}_{12}$ is
$m_1 \times n_2$, $\vec{A}_{21}$ is $m_2 \times n_1$, and
$\vec{A}_{22}$ is $m_2 \times n_2$, with $m_1 + m_2 = m$ and
$n_1 + n_2 = n$.

Two matrices $\vec{A}$ and $\vec{B}$ of the same dimensions that are
\emph{conformably partitioned} (i.e., partitioned so that corresponding
blocks have the same dimensions) can be added block by block:
\begin{equation}
  \vec{A} + \vec{B} =
  \begin{bmatrix}
    \vec{A}_{11} + \vec{B}_{11} & \vec{A}_{12} + \vec{B}_{12} \\
    \vec{A}_{21} + \vec{B}_{21} & \vec{A}_{22} + \vec{B}_{22}
  \end{bmatrix}.
  \label{eq:A.2}
\end{equation}

Partitioned matrices can also be multiplied block by block, provided that
the partitioning is \emph{conformable for multiplication}: the column
partition of the first matrix must match the row partition of the second.
If $\vec{A}$ is $m \times n$ and $\vec{B}$ is $n \times p$, with $\vec{A}$
partitioned as in \eqref{eq:A.1} and
\begin{equation}
  \vec{B} =
  \begin{bmatrix}
    \vec{B}_{11} & \vec{B}_{12} \\
    \vec{B}_{21} & \vec{B}_{22}
  \end{bmatrix},
  \label{eq:A.3}
\end{equation}
where $\vec{B}_{11}$ is $n_1 \times p_1$, $\vec{B}_{12}$ is
$n_1 \times p_2$, $\vec{B}_{21}$ is $n_2 \times p_1$, and
$\vec{B}_{22}$ is $n_2 \times p_2$, then the product $\vec{A}\vec{B}$
can be computed as
\begin{equation}
  \vec{A}\vec{B} =
  \begin{bmatrix}
    \vec{A}_{11}\vec{B}_{11} + \vec{A}_{12}\vec{B}_{21} &
      \vec{A}_{11}\vec{B}_{12} + \vec{A}_{12}\vec{B}_{22} \\
    \vec{A}_{21}\vec{B}_{11} + \vec{A}_{22}\vec{B}_{21} &
      \vec{A}_{21}\vec{B}_{12} + \vec{A}_{22}\vec{B}_{22}
  \end{bmatrix}.
  \label{eq:A.4}
\end{equation}
The key requirement is that the number of columns in $\vec{A}_{i1}$
equals the number of rows in $\vec{B}_{1j}$ (namely $n_1$), and the
number of columns in $\vec{A}_{i2}$ equals the number of rows in
$\vec{B}_{2j}$ (namely $n_2$).

As a special case, if $\vec{x}$ is an $n$-dimensional column vector
partitioned as
\[
  \vec{x} =
  \begin{bmatrix}
    \vec{x}_1 \\ \vec{x}_2
  \end{bmatrix},
  \quad
  \vec{x}_1 \; (n_1 \times 1), \quad
  \vec{x}_2 \; (n_2 \times 1),
\]
then
\begin{equation}
  \vec{A}\vec{x} =
  \begin{bmatrix}
    \vec{A}_{11}\vec{x}_1 + \vec{A}_{12}\vec{x}_2 \\
    \vec{A}_{21}\vec{x}_1 + \vec{A}_{22}\vec{x}_2
  \end{bmatrix}.
  \label{eq:A.5}
\end{equation}

%% ── A.2 Inverting Partitioned Matrices ───────────────────────────
\section{Inverting Partitioned Matrices}
\label{sec:A.2}

The inverse of a partitioned matrix can be expressed in terms of the
inverses of its blocks.  The key concept is the \emph{Schur complement}.

\begin{definition}[Schur Complement]
\label{def:schur}
Let
\[
  \vec{A} =
  \begin{bmatrix}
    \vec{A}_{11} & \vec{A}_{12} \\
    \vec{A}_{21} & \vec{A}_{22}
  \end{bmatrix}
\]
be a nonsingular matrix with $\vec{A}_{11}$ and $\vec{A}_{22}$ square.
If $\vec{A}_{22}$ is nonsingular, the \emph{Schur complement of
$\vec{A}_{22}$ in $\vec{A}$} is defined as
\begin{equation}
  \vec{A}_{11 \cdot 2}
  \equiv \vec{A}_{11} - \vec{A}_{12}\vec{A}_{22}^{-1}\vec{A}_{21}.
  \label{eq:A.6}
\end{equation}
Similarly, if $\vec{A}_{11}$ is nonsingular, the \emph{Schur complement
of $\vec{A}_{11}$ in $\vec{A}$} is
\begin{equation}
  \vec{A}_{22 \cdot 1}
  \equiv \vec{A}_{22} - \vec{A}_{21}\vec{A}_{11}^{-1}\vec{A}_{12}.
  \label{eq:A.7}
\end{equation}
\end{definition}

\begin{proposition}[Partitioned Inverse]
\label{prop:partitioned-inverse}
If $\vec{A}$ and $\vec{A}_{22}$ are both nonsingular, then
\begin{equation}
  \vec{A}^{-1} =
  \begin{bmatrix}
    \vec{A}_{11 \cdot 2}^{-1} &
      -\vec{A}_{11 \cdot 2}^{-1}\vec{A}_{12}\vec{A}_{22}^{-1} \\[4pt]
    -\vec{A}_{22}^{-1}\vec{A}_{21}\vec{A}_{11 \cdot 2}^{-1} &
      \vec{A}_{22}^{-1} + \vec{A}_{22}^{-1}\vec{A}_{21}
        \vec{A}_{11 \cdot 2}^{-1}\vec{A}_{12}\vec{A}_{22}^{-1}
  \end{bmatrix}.
  \label{eq:A.8}
\end{equation}
Alternatively, if $\vec{A}$ and $\vec{A}_{11}$ are both nonsingular, then
\begin{equation}
  \vec{A}^{-1} =
  \begin{bmatrix}
    \vec{A}_{11}^{-1} + \vec{A}_{11}^{-1}\vec{A}_{12}
      \vec{A}_{22 \cdot 1}^{-1}\vec{A}_{21}\vec{A}_{11}^{-1} &
      -\vec{A}_{11}^{-1}\vec{A}_{12}\vec{A}_{22 \cdot 1}^{-1} \\[4pt]
    -\vec{A}_{22 \cdot 1}^{-1}\vec{A}_{21}\vec{A}_{11}^{-1} &
      \vec{A}_{22 \cdot 1}^{-1}
  \end{bmatrix}.
  \label{eq:A.9}
\end{equation}
\end{proposition}

\begin{proof}
To verify \eqref{eq:A.8}, multiply the right-hand side by $\vec{A}$ and
confirm the product is the identity matrix.  Write
\[
  \begin{bmatrix}
    \vec{A}_{11 \cdot 2}^{-1} &
      -\vec{A}_{11 \cdot 2}^{-1}\vec{A}_{12}\vec{A}_{22}^{-1} \\[4pt]
    -\vec{A}_{22}^{-1}\vec{A}_{21}\vec{A}_{11 \cdot 2}^{-1} &
      \vec{A}_{22}^{-1} + \vec{A}_{22}^{-1}\vec{A}_{21}
        \vec{A}_{11 \cdot 2}^{-1}\vec{A}_{12}\vec{A}_{22}^{-1}
  \end{bmatrix}
  \begin{bmatrix}
    \vec{A}_{11} & \vec{A}_{12} \\
    \vec{A}_{21} & \vec{A}_{22}
  \end{bmatrix}.
\]
For the $(1,1)$ block of the product:
\begin{align*}
  &\vec{A}_{11 \cdot 2}^{-1}\vec{A}_{11}
    - \vec{A}_{11 \cdot 2}^{-1}\vec{A}_{12}\vec{A}_{22}^{-1}\vec{A}_{21} \\
  &\quad = \vec{A}_{11 \cdot 2}^{-1}
    (\vec{A}_{11} - \vec{A}_{12}\vec{A}_{22}^{-1}\vec{A}_{21})
  = \vec{A}_{11 \cdot 2}^{-1}\vec{A}_{11 \cdot 2}
  = \vec{I}.
\end{align*}
For the $(1,2)$ block:
\begin{align*}
  &\vec{A}_{11 \cdot 2}^{-1}\vec{A}_{12}
    - \vec{A}_{11 \cdot 2}^{-1}\vec{A}_{12}\vec{A}_{22}^{-1}\vec{A}_{22}
  = \vec{A}_{11 \cdot 2}^{-1}\vec{A}_{12}
    - \vec{A}_{11 \cdot 2}^{-1}\vec{A}_{12}
  = \vec{0}.
\end{align*}
The $(2,1)$ and $(2,2)$ blocks can be verified similarly.
Formula~\eqref{eq:A.9} is verified in the same manner.
\end{proof}

By comparing the $(1,1)$ blocks of \eqref{eq:A.8} and \eqref{eq:A.9},
one obtains a useful identity.

\begin{proposition}[Matrix Inversion Lemma]
\label{prop:matrix-inversion-lemma}
If $\vec{A}_{11}$, $\vec{A}_{22}$, and $\vec{A}$ are all nonsingular,
then
\begin{equation}
  (\vec{A}_{11} - \vec{A}_{12}\vec{A}_{22}^{-1}\vec{A}_{21})^{-1}
  = \vec{A}_{11}^{-1} + \vec{A}_{11}^{-1}\vec{A}_{12}
    (\vec{A}_{22} - \vec{A}_{21}\vec{A}_{11}^{-1}\vec{A}_{12})^{-1}
    \vec{A}_{21}\vec{A}_{11}^{-1}.
  \label{eq:A.10}
\end{equation}
That is,
\begin{equation}
  \vec{A}_{11 \cdot 2}^{-1}
  = \vec{A}_{11}^{-1} + \vec{A}_{11}^{-1}\vec{A}_{12}\,
    \vec{A}_{22 \cdot 1}^{-1}\,
    \vec{A}_{21}\vec{A}_{11}^{-1}.
  \label{eq:A.11}
\end{equation}
\end{proposition}

This result is sometimes called the \emph{Woodbury identity} or the
\emph{Sherman--Morrison--Woodbury formula}.  It is especially useful
when $\vec{A}_{11}$ is large and easy to invert while
$\vec{A}_{22}$ is small.

%% ── A.3 Kronecker Products ──────────────────────────────────────
\section{Kronecker Products}
\label{sec:A.3}

\begin{definition}[Kronecker Product]
\label{def:kronecker}
Let $\vec{A}$ be $m \times n$ and $\vec{B}$ be $p \times q$.
The \emph{Kronecker product} (or \emph{direct product})
$\vec{A} \otimes \vec{B}$ is the $mp \times nq$ matrix defined by
\begin{equation}
  \vec{A} \otimes \vec{B}
  =
  \begin{bmatrix}
    a_{11}\vec{B} & a_{12}\vec{B} & \cdots & a_{1n}\vec{B} \\
    a_{21}\vec{B} & a_{22}\vec{B} & \cdots & a_{2n}\vec{B} \\
    \vdots & \vdots & \ddots & \vdots \\
    a_{m1}\vec{B} & a_{m2}\vec{B} & \cdots & a_{mn}\vec{B}
  \end{bmatrix},
  \label{eq:A.12}
\end{equation}
where $a_{ij}$ is the $(i,j)$ element of $\vec{A}$.
\end{definition}

The Kronecker product has several useful properties.

\begin{proposition}[Properties of Kronecker Products]
\label{prop:kronecker-properties}
Let $\vec{A}$, $\vec{B}$, $\vec{C}$, and $\vec{D}$ be matrices whose
dimensions are such that the indicated operations are well defined, and
let $c$ be a scalar.  Then:
\begin{enumerate}[label=(\alph*)]
  \item \textbf{Associativity}:
    \begin{equation}
      (\vec{A} \otimes \vec{B}) \otimes \vec{C}
      = \vec{A} \otimes (\vec{B} \otimes \vec{C}).
      \label{eq:A.13}
    \end{equation}

  \item \textbf{Distributivity over addition}:
    \begin{equation}
      \vec{A} \otimes (\vec{B} + \vec{C})
      = \vec{A} \otimes \vec{B} + \vec{A} \otimes \vec{C},
      \label{eq:A.14}
    \end{equation}
    \begin{equation}
      (\vec{A} + \vec{B}) \otimes \vec{C}
      = \vec{A} \otimes \vec{C} + \vec{B} \otimes \vec{C}.
      \label{eq:A.15}
    \end{equation}

  \item \textbf{Scalar multiplication}:
    \begin{equation}
      c(\vec{A} \otimes \vec{B})
      = (c\vec{A}) \otimes \vec{B}
      = \vec{A} \otimes (c\vec{B}).
      \label{eq:A.16}
    \end{equation}

  \item \textbf{Mixed-product property}: If $\vec{A}$ is $m \times n$,
    $\vec{B}$ is $p \times q$, $\vec{C}$ is $n \times r$, and
    $\vec{D}$ is $q \times s$, then
    \begin{equation}
      (\vec{A} \otimes \vec{B})(\vec{C} \otimes \vec{D})
      = (\vec{A}\vec{C}) \otimes (\vec{B}\vec{D}).
      \label{eq:A.17}
    \end{equation}

  \item \textbf{Transpose}:
    \begin{equation}
      (\vec{A} \otimes \vec{B})'
      = \vec{A}' \otimes \vec{B}'.
      \label{eq:A.18}
    \end{equation}

  \item \textbf{Inverse}: If $\vec{A}$ ($m \times m$) and $\vec{B}$
    ($p \times p$) are both nonsingular, then $\vec{A} \otimes \vec{B}$
    is nonsingular and
    \begin{equation}
      (\vec{A} \otimes \vec{B})^{-1}
      = \vec{A}^{-1} \otimes \vec{B}^{-1}.
      \label{eq:A.19}
    \end{equation}

  \item \textbf{Trace}: If $\vec{A}$ is $m \times m$ and $\vec{B}$ is
    $p \times p$, then
    \begin{equation}
      \tr(\vec{A} \otimes \vec{B})
      = \tr(\vec{A}) \cdot \tr(\vec{B}).
      \label{eq:A.20}
    \end{equation}

  \item \textbf{Determinant}: If $\vec{A}$ is $m \times m$ and
    $\vec{B}$ is $p \times p$, then
    \begin{equation}
      \det(\vec{A} \otimes \vec{B})
      = (\det \vec{A})^{p} \cdot (\det \vec{B})^{m}.
      \label{eq:A.21}
    \end{equation}

  \item \textbf{Rank}:
    \begin{equation}
      \rank(\vec{A} \otimes \vec{B})
      = \rank(\vec{A}) \cdot \rank(\vec{B}).
      \label{eq:A.22}
    \end{equation}
\end{enumerate}
\end{proposition}

\begin{proof}
Properties (a)--(c) follow directly from the definition~\eqref{eq:A.12}.

\medskip
\noindent\textbf{(d) Mixed-product property.}\quad
The $(i,j)$ block of $(\vec{A} \otimes \vec{B})(\vec{C} \otimes \vec{D})$
is
\[
  \sum_{k=1}^{n} a_{ik}\vec{B} \cdot c_{kj}\vec{D}
  = \biggl(\sum_{k=1}^{n} a_{ik}c_{kj}\biggr) \vec{B}\vec{D}
  = [\vec{A}\vec{C}]_{ij}\,(\vec{B}\vec{D}),
\]
which is exactly the $(i,j)$ block of
$(\vec{A}\vec{C}) \otimes (\vec{B}\vec{D})$.

\medskip
\noindent\textbf{(e) Transpose.}\quad
Follows from the definition, since the $(i,j)$ block of
$(\vec{A} \otimes \vec{B})'$ is $(a_{ji}\vec{B})' = a_{ji}\vec{B}'$,
which equals the $(i,j)$ block of $\vec{A}' \otimes \vec{B}'$.

\medskip
\noindent\textbf{(f) Inverse.}\quad
By the mixed-product property~(d),
\[
  (\vec{A} \otimes \vec{B})(\vec{A}^{-1} \otimes \vec{B}^{-1})
  = (\vec{A}\vec{A}^{-1}) \otimes (\vec{B}\vec{B}^{-1})
  = \vec{I}_m \otimes \vec{I}_p
  = \vec{I}_{mp}.
\]

\medskip
\noindent\textbf{(g) Trace.}\quad
Since $\tr(\vec{A} \otimes \vec{B})
= \sum_{i=1}^{m} \tr(a_{ii}\vec{B})
= \sum_{i=1}^{m} a_{ii}\,\tr(\vec{B})
= \tr(\vec{A})\,\tr(\vec{B})$.

\medskip
\noindent\textbf{(h) Determinant.}\quad
See, e.g., Horn and Johnson (1985, Theorem~4.2.15).

\medskip
\noindent\textbf{(i) Rank.}\quad
See, e.g., Horn and Johnson (1985, Corollary~4.2.11).
\end{proof}

A useful identity relating the Kronecker product to the $\vect$
operator (which stacks the columns of a matrix into a single column
vector) is the following.

\begin{proposition}[Vec and Kronecker Product]
\label{prop:vec-kronecker}
For matrices $\vec{A}$ ($m \times n$), $\vec{B}$ ($n \times p$), and
$\vec{C}$ ($p \times q$),
\begin{equation}
  \vect(\vec{A}\vec{B}\vec{C})
  = (\vec{C}' \otimes \vec{A})\,\vect(\vec{B}).
  \label{eq:A.23}
\end{equation}
As a special case, for any $n \times 1$ vector $\vec{b}$,
\begin{equation}
  \vec{A}\vec{b} = (\vec{b}' \otimes \vec{I}_m)
    \,\vect(\vec{A})
  = (\vec{I}_n \otimes \vec{A})\,\vec{b}.
  \label{eq:A.24}
\end{equation}
\end{proposition}

\begin{proof}
Let $\vec{c}_j$ denote the $j$-th column of $\vec{C}$.  The $j$-th
column of $\vec{A}\vec{B}\vec{C}$ is $\vec{A}\vec{B}\vec{c}_j$.
Therefore
\[
  \vect(\vec{A}\vec{B}\vec{C})
  = \begin{bmatrix}
    \vec{A}\vec{B}\vec{c}_1 \\
    \vec{A}\vec{B}\vec{c}_2 \\
    \vdots \\
    \vec{A}\vec{B}\vec{c}_q
  \end{bmatrix}
  = \begin{bmatrix}
    \vec{c}_1' \otimes \vec{A} \\
    \vec{c}_2' \otimes \vec{A} \\
    \vdots \\
    \vec{c}_q' \otimes \vec{A}
  \end{bmatrix}
  \vect(\vec{B})
  = (\vec{C}' \otimes \vec{A})\,\vect(\vec{B}). \qedhere
\]
\end{proof}
